\subsubsection{Quaternions e justificativa de sua utilização}
\label{sec:representacao_quaternion}
 
Nessas condições, os quaternions constituem uma alternativa conveniente para representar a orientação, pois evitam as descontinuidades associadas às singularidades dos ângulos de Euler e permitem uma descrição globalmente válida da atitude. Um quaternion unitário é definido como um elemento do espaço quaternional de quatro dimensões, pertencente ao conjunto
dos quaternions unitários $q = [w,\vec{v}\,]$, com $\|q\| = 1$, em que $w$ é a parte escalar e $\vec{v}$ é a parte vetorial. Cada quaternion unitário $q$ corresponde a uma rotação de ângulo $\theta$ em torno de um eixo unitário $\vec{u}$, dada por
\begin{equation}
q(\theta,\vec{u}) = \cos\!\left(\frac{\theta}{2}\right) + \sin\!\left(\frac{\theta}{2}\right)\,\vec{u}.
\label{eq:quat_axis_angle}
\end{equation}

Essa parametrização fornece uma descrição global da atitude, livre de singularidades geométricas e numericamente estável, adequada para a integração contínua das equações diferenciais de atitude. A normalização unitária $\|q\|=1$ elimina redundâncias e permite utilizar a própria condição de norma como verificação da qualidade numérica da integração ao longo de simulações prolongadas.

Do ponto de vista físico, os quaternions preservam a norma e realizam rotações por multiplicação, assegurando uma transição suave entre orientações sucessivas. Embora exijam uma variável adicional em comparação aos ângulos de Euler, sua principal vantagem reside na ausência de singularidades e no fato de que as equações cinemáticas quaternionais envolvem apenas produtos, sem funções trigonométricas, o que é particularmente conveniente para implementação embarcada em tempo real \cite{wie2008space}. Além disso, operações como composição de rotações, inversão e interpolação são tratadas de forma direta na álgebra quaternional \cite{gonzalez2025quaternionic}.

Por essas razões, a formulação cinemática e dinâmica de atitude adotada neste trabalho utiliza quaternions unitários, tanto para a base da espaçonave perseguidora quanto para as transformações locais das juntas do \gls{mr}, evitando ambiguidades associadas a parametrizações angulares e garantindo continuidade global na descrição do movimento.

\subsubsection{Definição e componentes do quaternion}
\label{sec:componentes_quaternion}

O espaço quaternional é um espaço vetorial de quatro dimensões, composto por um elemento escalar e três elementos imaginários mutuamente ortogonais $\{i,j,k\}$, que obedecem às regras de multiplicação introduzidas por Hamilton (1844):
\begin{equation}
i^2 = j^2 = k^2 = ijk = -1,
\label{eq:hamilton_quadrados}
\end{equation}
\begin{equation}
ij = k,\quad jk = i,\quad ki = j,\quad
ji = -k,\quad kj = -i,\quad ik = -j.
\label{eq:hamilton_produtos}
\end{equation}

Um quaternion genérico pode ser escrito como
\begin{equation}
q = w + x\,i + y\,j + z\,k,
\label{eq:quaternion_forma}
\end{equation}
ou, de forma vetorial compacta,
\begin{equation}
q = [\,w,\vec{v}\,] = [\,w,x,y,z\,],
\label{eq:quaternion_wxyz}
\end{equation}
em que $w$ é a parte escalar e $\vec{v} = [x,y,z]$ é a parte vetorial. No contexto deste trabalho, será adotada explicitamente a notação $q = [w,x,y,z]^T$, em conformidade com a convenção usual em aplicações aeroespaciais e robóticas.

A magnitude (ou norma) de um quaternion é definida por
\begin{equation}
\|q\| = \sqrt{w^2 + x^2 + y^2 + z^2},
\label{eq:norma_quaternion}
\end{equation}
e o quaternion unitário é obtido normalizando-se $q$ de modo que $\|q\| = 1$. Em aplicações de atitude, essa condição de unidade é preservada pela dinâmica exata e pode ser utilizada como critério de monitoramento de erros numéricos.

A relação entre os quaternions e a velocidade angular do corpo é obtida a partir da equação cinemática diferencial \cite{wie2008space}. Considerando o vetor de estado quaternional
\begin{equation}
\mathbf{q} =
\begin{bmatrix}
w \\ x \\ y \\ z
\end{bmatrix},
\qquad
\vec{\omega} =
\begin{bmatrix}
\omega_x \\ \omega_y \\ \omega_z
\end{bmatrix},
\label{eq:q_omega_vetores}
\end{equation}
a equação cinemática pode ser escrita na forma matricial
\begin{equation}
\begin{bmatrix}
\dot{w} \\[2pt]
\dot{x} \\[2pt]
\dot{y} \\[2pt]
\dot{z}
\end{bmatrix}
=
\frac{1}{2}
\begin{bmatrix}
0 & -\omega_x & -\omega_y & -\omega_z \\
\omega_x & 0 & \omega_z & -\omega_y \\
\omega_y & -\omega_z & 0 & \omega_x \\
\omega_z & \omega_y & -\omega_x & 0
\end{bmatrix}
\begin{bmatrix}
w \\ x \\ y \\ z
\end{bmatrix},
\label{eq:cinematica_quaternion_wxyz}
\end{equation}
a qual é equivalente à forma clássica apresentada por Wie em termos dos componentes $q_1,q_2,q_3,q_4$ do quaternion \cite{wie2008space}.

De forma compacta, pode-se escrever
\begin{equation}
\dot{\mathbf{q}} = \tfrac{1}{2}\,\boldsymbol{\Omega}(\vec{\omega})\,\mathbf{q},
\label{eq:quat_compacta}
\end{equation}
em que $\boldsymbol{\Omega}(\vec{\omega})$ é a matriz $4\times 4$ definida em \eqref{eq:cinematica_quaternion_wxyz}. Essa estrutura garante que, se $\|q(0)\|=1$, a solução exata de \eqref{eq:quat_compacta} preserva a unidade para todo $t>0$, propriedade que é explorada como verificação de consistência numérica nas integrações de atitude.
